% Provenance (never printed): sources/book2/BRIEFS.md, g12-15. Life satisfaction against income across countries and inside
% them; the log relation (proved: the same gain per doubling); what GDP leaves out, with numbers (free time, health, the
% distribution, nature); other indicators and how closely they follow income; what economics measures and what it does not.
% Real numbers: figdata/grade-12/15-income-and-wellbeing/ (Eurostat ilc_pw01 life satisfaction 2018 with WDI GDP per person
% at purchasing power; Penn World Table hours; WDI life expectancy; five invented countries, seeded). Ledger: Stevenson and
% Wolfers 2008, Deaton 2008, Killingsworth 2021, Killingsworth-Kahneman-Mellers 2023 (abstracts), Easterlin 1974 (Wikipedia,
% dated label). No photograph; no debate box (Year 3 owns the open question).
\chapter{Income and Wellbeing}\label{ch:g12:income-and-wellbeing}

In 2018 people in Finland, asked to rate their satisfaction with their life from 0 to 10, gave on average 8.1. In Bulgaria they gave 5.4. Finland produced about 56,000
dollars a person, at purchasing power, and Bulgaria about 28,000. Across 37 European countries, the richer the country, the more satisfied its people said they were.
% figdata: well_key (fi_s 8.1, fi_y 56226, bg_s 5.4, bg_y 27691, n_ls 37)

This last chapter asks how closely income follows what people value, and what the measures of this book leave out.

\begin{openingchart}{Mean life satisfaction (0--10) against GDP per person at constant purchasing power (logarithmic scale), 37 European countries, 2018;
the line is the fit $s = a + b \ln y$.}{Eurostat (ilc\_pw01); World Bank, World Development Indicators (NY.GDP.PCAP.PP.KD)}
\begin{tikzpicture}
\begin{axis}[omscatter, width=0.8\linewidth, height=6.0cm, xmode=log, xmin=9000, xmax=150000, ymin=5, ymax=8.5, log ticks with fixed point,
  xtick={10000,20000,50000,100000}, xticklabels={10, 20, 50, 100}, xlabel={GDP per person, thousands of dollars}, ylabel={life satisfaction}]
\addplot+[omDef, mark=*, mark size=1.3pt, only marks] table[x=gdppc, y=lifesat] {figdata/grade-12/15-income-and-wellbeing/out/lifesat.csv};
\addplot[black, dashed, sharp plot, no marks] table[x=gdppc, y=lifesat] {figdata/grade-12/15-income-and-wellbeing/out/fit.csv};
\node[font=\scriptsize, anchor=north] at (axis cs:27691,5.4) {Bulgaria};
\node[font=\scriptsize, anchor=south east] at (axis cs:34829,7.4) {Romania};
\node[font=\scriptsize, anchor=south east] at (axis cs:56226,8.1) {Finland};
\node[font=\scriptsize, anchor=north east] at (axis cs:131734,7.45) {Luxembourg};
\node[font=\scriptsize, anchor=south] at (axis cs:11188,5.9) {Kosovo};
\end{axis}
\end{tikzpicture}
\end{openingchart}

\begin{recall}
From the grades 6--9 volume: what GDP leaves out, and life expectancy. From \cref{ch:g10:national-accounts}: how GDP is measured. From \cref{ch:g11:inequality}: the \omterm{def:g11:inequality:gini}{Gini
coefficient}. From mathematics: the logarithm.
\end{recall}

\section{Asking people}\label{sec:g12:income-and-wellbeing:asking}

\begin{definition}[Life satisfaction]\label{def:g12:income-and-wellbeing:life-satisfaction}
\emph{Life satisfaction}\index{life satisfaction} is a person's own rating of their life as a whole, given in answer to a survey question on a fixed scale.
\end{definition}

\begin{definition}[Subjective wellbeing]\label{def:g12:income-and-wellbeing:subjective-wellbeing}
\emph{Subjective wellbeing}\index{subjective wellbeing} is wellbeing as people report it themselves: their \omterm{def:g12:income-and-wellbeing:life-satisfaction}{life satisfaction}, and the feelings, such as happiness, worry or
stress, they report over a day.
\end{definition}

\begin{measure}[How life satisfaction is measured]
A large sample of people, chosen at random, are asked to rate their satisfaction with their life overall on a scale from 0 (not at all satisfied) to 10 (completely
satisfied); the average is published by country and by group. Another method asks people at random moments, often by phone, how they feel at that moment. The answers depend on
how the question is worded and on where it comes in the survey, so comparisons are safest when the same question is asked in the same way; whether people in different
countries use the scale in the same way cannot be checked directly.
\end{measure}

The relation in the opening chart is not only between countries. Within countries, richer people report higher satisfaction than poorer people, and the slope is similar to the
slope across countries.
% ledger: g12:sw2008

\section{A relation in the logarithm}\label{sec:g12:income-and-wellbeing:log}

\begin{proposition}[The same gain for each doubling]\label{prop:g12:income-and-wellbeing:doubling}
If \omterm{def:g12:income-and-wellbeing:life-satisfaction}{life satisfaction} $s$ and income $y$ satisfy $s = a + b \ln y$, then multiplying income by any factor $k$ raises $s$ by $b \ln k$, whatever the starting income. In
particular, each doubling adds $b \ln 2$.
\end{proposition}
\begin{proof}
$s(ky) - s(y) = a + b \ln k + b \ln y - a - b \ln y = b \ln k$.
\end{proof}

On the opening chart, $b = 1.35$: each doubling of income goes with a rise of 0.93 points. A rise of 1,000 dollars therefore matters much more in a poor country than in a rich
one, because it is a larger fraction of income, while a 10\,\% rise matters as much everywhere.
% figdata: well_key (b_ln 1.346, per_doubling 0.93, corr_ls 0.84)

\begin{omfigure}
\begin{tikzpicture}
\begin{axis}[name=lin, width=0.46\linewidth, height=4.6cm, xmin=0, xmax=64, ymin=4, ymax=8, xlabel={income, thousands}, ylabel={satisfaction}, xtick={4,16,32,64},
  title={\scriptsize linear axis}, grid=major, grid style={gray!20}]
\addplot[omDef, thick, domain=2:64, samples=80] {4.5 + 0.6*ln(x/2)/ln(2)};
\end{axis}
\begin{axis}[at={(lin.east)}, anchor=west, xshift=1.2cm, width=0.46\linewidth, height=4.6cm, xmode=log, log basis x=2, xmin=2, xmax=64, ymin=4, ymax=8,
  xlabel={income, thousands}, xtick={2,4,8,16,32,64}, xticklabels={2,4,8,16,32,64}, title={\scriptsize logarithmic axis}, grid=major, grid style={gray!20}]
\addplot[omDef, thick, domain=2:64, samples=2] {4.5 + 0.6*ln(x/2)/ln(2)};
\end{axis}
\end{tikzpicture}

{\small The same relation, $s = 4.5 + 0.6 \log_2(y/2)$, on two axes. On a linear axis it flattens as income rises; on a logarithmic axis, where each doubling takes the same
width, it is a straight line: 0.6 points per doubling, at 2,000 as at 32,000.}
\end{omfigure}

\begin{model}{Life satisfaction in the logarithm of income}
\assumes Average \omterm{def:g12:income-and-wellbeing:life-satisfaction}{life satisfaction} rises with the logarithm of income.
\predicts The same gain for each doubling of income, at any income, across countries and within them.
\tested Consistent with it: across countries, a doubling of income goes with nearly one more point on the 0--10 scale, a slope no smaller among rich countries; no point has been found beyond
  which it stops. In the United States, reports of how people feel during the day rise in a straight line with the logarithm of income, even above 75,000 dollars a year.
\breaks For the least happy people, whose good feelings stop rising with income at high incomes. And perhaps over time: a 1974 study found that within a
  country, average happiness did not trend upward as income grew (Easterlin 1974); later work with other data found a positive relation over time, and the
  question is still argued over.
\end{model}
% ledger: g12:deaton2008, g12:sw2008, g12:kill2021, g12:kkm2023, g12:easterlin1974

A slope across countries is not the effect of raising a country's income. Richer countries also have better health, more trust and more security, which people may value for
themselves; and satisfied, healthy people may earn more. The chart shows how closely the two move together, not how much one causes the other.

\section{What GDP leaves out}\label{sec:g12:income-and-wellbeing:leaves-out}

\emph{Free time.} GDP counts what is produced, not the hours it took. In Germany the average worker worked 2,427 hours in 1950 and 1,372 in 2019, while output per person rose
7.2 times; in the United States, hours fell from 2,022 to 1,797 while output per person rose 3.8 times. Part of the gain from growth was taken as time. Across countries in 2019
the link is weak (a correlation of $-0.18$ over 130 countries): in the Penn World Table's estimates, Mexicans worked fewer hours than Koreans, though Koreans were more
than twice as rich. Other sources measure hours differently, and such comparisons across countries are rough.
% figdata: well_key (h1950_DEU 2427, h2019_DEU 1372, ymult_DEU 7.2, h1950_USA 2022, h2019_USA 1797, ymult_USA 3.8, n_hr 130, corr_hr -0.18); hours.csv (MEX 1639 at 21299, KOR 1963 at 49200)

\begin{omchart}{Average annual hours worked per person employed, six economies, 1950--2019 (South Korea from 1953). Straight stretches in the early years, such as
Germany's before 1960 and South Korea's before 1963, are interpolations between benchmark years in the source, not yearly data.}{Penn World Table 11.0 (avh)}
\begin{tikzpicture}
\begin{axis}[omtime, width=0.84\linewidth, height=5.6cm, xmin=1950, xmax=2019, ymin=1200, ymax=3100, ytick={1200,1600,2000,2400,2800}, xtick={1950,1960,...,2010}, ylabel={hours a year},
  y tick label style={/pgf/number format/.cd, 1000 sep={}}, unbounded coords=discard]
\addplot[omMeth, thick] table[x=year, y=DEU] {figdata/grade-12/15-income-and-wellbeing/out/hours_ts.csv};
\addplot[gray, thick] table[x=year, y=FRA] {figdata/grade-12/15-income-and-wellbeing/out/hours_ts.csv};
\addplot[black, thick, dashed] table[x=year, y=GBR] {figdata/grade-12/15-income-and-wellbeing/out/hours_ts.csv};
\addplot[omProp, thick] table[x=year, y=JPN] {figdata/grade-12/15-income-and-wellbeing/out/hours_ts.csv};
\addplot[omThm, thick] table[x=year, y=KOR] {figdata/grade-12/15-income-and-wellbeing/out/hours_ts.csv};
\addplot[omDef, thick] table[x=year, y=USA] {figdata/grade-12/15-income-and-wellbeing/out/hours_ts.csv};
\node[font=\scriptsize, omThm, anchor=south] at (axis cs:1975,2930) {South Korea};
\node[font=\scriptsize, omProp, anchor=south west] at (axis cs:1960,2240) {Japan};
\draw[gray] (axis cs:1951.5,2410) -- (axis cs:1951.5,2840) node[anchor=south west, xshift=-2pt, font=\scriptsize, omMeth] {Germany};
\draw[gray] (axis cs:1956,1985) -- (axis cs:1956,1720) node[anchor=north, font=\scriptsize, omDef] {United States};
\node[font=\scriptsize, anchor=north] at (axis cs:1990,1420) {France (grey), United Kingdom (dashed)};
\end{axis}
\end{tikzpicture}
\end{omchart}

\emph{Health.} Life expectancy is not in GDP, but it follows income closely: across 198 countries in 2019 the correlation between life expectancy and the logarithm of income
per person was 0.86. The chart below shows the cloud: it rises steeply among the poorest countries and flattens among the richest,
even on a logarithmic scale. Countries differ around it. Costa Rica, with a third of American income, had a longer life expectancy, 80.3 years against 78.8; Nigeria, with
a little more income than India, had 53.0 years against 70.7. People's own ratings, though, respond little to measured health: neither \omterm{def:g12:income-and-wellbeing:life-satisfaction}{life satisfaction} nor
satisfaction with health responds strongly to objective measures of health.
% figdata: well_key (n_le 198, corr_le 0.86); lifeexp (CRI 23497 80.3; USA 69121 78.8; NGA 8362 53.0; IND 7702 70.7); ledger: g12:deaton2008

\begin{omchart}{Life expectancy at birth against GDP per person at constant purchasing power (logarithmic scale), 198 countries, 2019. The Central African Republic
is left out: its series jumps by twenty years from one year to the next.}{World Bank, World Development Indicators (SP.DYN.LE00.IN, NY.GDP.PCAP.PP.KD)}
\begin{tikzpicture}
\begin{axis}[omscatter, width=0.8\linewidth, height=6.0cm, xmode=log, xmin=800, xmax=160000, ymin=50, ymax=88, log ticks with fixed point,
  xtick={1000,3000,10000,30000,100000}, xticklabels={1, 3, 10, 30, 100}, xlabel={GDP per person, thousands of dollars}, ylabel={years}]
\addplot+[omDef, mark=*, mark size=1.1pt, only marks] table[x=gdppc, y=le] {figdata/grade-12/15-income-and-wellbeing/out/lifeexp.csv};
\draw[gray] (axis cs:69121,78.8) -- (axis cs:100000,70) node[anchor=north, font=\scriptsize, black] {United States};
\draw[gray] (axis cs:23497,80.3) -- (axis cs:9000,85) node[anchor=east, font=\scriptsize, black] {Costa Rica};
\draw[gray] (axis cs:46055,84.4) -- (axis cs:30000,87) node[anchor=east, font=\scriptsize, black] {Japan};
\draw[gray] (axis cs:8362,53.0) -- (axis cs:20000,54) node[anchor=west, font=\scriptsize, black] {Nigeria};
\draw[gray] (axis cs:7702,70.7) -- (axis cs:2500,77) node[anchor=east, font=\scriptsize, black] {India};
\end{axis}
\end{tikzpicture}
\end{omchart}

\emph{The distribution.} GDP per person is an average. Two countries with the same average can differ in their \omterm{def:g11:inequality:gini}{Gini coefficient}, and if satisfaction rises with the logarithm
of income, a dollar moved from a rich person to a poor one raises the average log of income: the same average income goes with lower average satisfaction where it is
spread more unequally. Take two invented countries of five people with the same average income of 20,000 dollars: in one each person has 20,000, in the other
they have 5, 10, 15, 25 and 45 thousand. The second country's average log income is that of 15,319 dollars, its geometric mean, and with the opening chart's slope of
1.35 its average satisfaction is $1.35 \ln(20\,000/15\,319) \approx 0.36$ points lower. Moving income has costs of its own, the \omterm{def:g10:what-prices-do:dwl}{deadweight loss} of the taxes that pay for
it (\cref{ch:g10:tax-incidence}) and the weaker reason to work and save, which this comparison leaves out.
% figdata: dist (geo_mean 15319, sat_gap 0.36); well_key (b_ln 1.346)

\emph{Nature.} GDP counts the timber sold, not the forest lost; it counts the output of a factory, not the harm its smoke does. Carbon dioxide is the largest such
harm by most estimates, and it rises with income: in 2019 each American emitted 15.6 tonnes from fuels and industry, each Swede 3.5 and each Indian 1.8. Across 152
countries of a million people or more, emissions per person rose slightly faster than income, by about 1.2\,\% for each 1\,\% of income, with a correlation of 0.91 in
logarithms. The World Bank values the damage at 40 dollars a tonne: on that price, it took 0.9\,\% of national income in the United States in 2019, 3.1\,\% in China,
4.2\,\% in Russia and 4.7\,\% in South Africa, whose output uses more coal and oil per dollar. The price is an estimate, and a different one would scale these numbers;
earlier chapters showed how a tax or a permit market can put such a price on the harm itself (\cref{ch:g11:externalities}). Other costs, such as a lost species, have no
agreed price at all.
% figdata: co2 (USA 15.61, SWE 3.53, IND 1.79); nature_key (n_co2 152, slope_co2 1.18, corr_co2 0.91); composite (dmg USA 0.9, CHN 3.1, RUS 4.2, ZAF 4.7)

\begin{omchart}{Carbon dioxide emitted per person against GDP per person at constant purchasing power, both on logarithmic scales, 152 countries of a million
people or more, 2019.}{Community Emissions Data System (CEDS), carbon dioxide by country; World Bank, World Development Indicators (SP.POP.TOTL, NY.GDP.PCAP.PP.KD)}
\begin{tikzpicture}
\begin{axis}[omscatter, width=0.8\linewidth, height=6.0cm, xmode=log, ymode=log, xmin=800, xmax=160000, ymin=0.02, ymax=60, log ticks with fixed point,
  xtick={1000,3000,10000,30000,100000}, xticklabels={1, 3, 10, 30, 100}, ytick={0.1,1,10}, xlabel={GDP per person, thousands of dollars},
  ylabel={tonnes per person}]
\addplot+[omDef, mark=*, mark size=1.1pt, only marks] table[x=gdppc, y=t] {figdata/grade-12/15-income-and-wellbeing/out/co2.csv};
\draw[gray] (axis cs:69121,15.61) -- (axis cs:30000,35) node[anchor=east, font=\scriptsize, black] {United States};
\draw[gray] (axis cs:61071,3.53) -- (axis cs:100000,1) node[anchor=north, font=\scriptsize, black] {Sweden};
\draw[gray] (axis cs:23497,1.63) -- (axis cs:30000,0.3) node[anchor=north, font=\scriptsize, black] {Costa Rica};
\draw[gray] (axis cs:7702,1.79) -- (axis cs:2500,8) node[anchor=east, font=\scriptsize, black] {India};
\draw[gray] (axis cs:18821,7.77) -- (axis cs:6000,25) node[anchor=east, font=\scriptsize, black] {China};
\end{axis}
\end{tikzpicture}
\end{omchart}

\section{Other indicators}\label{sec:g12:income-and-wellbeing:indicators}

\begin{definition}[Composite indicator]\label{def:g12:income-and-wellbeing:composite-indicator}
A \emph{composite indicator}\index{composite indicator} combines several measures, each rescaled to a common range, into a single number by a weighted average.
\end{definition}

Each measure is rescaled between two bounds, a minimum and a maximum, so that the index
\[ \frac{x - x_{\min}}{x_{\max} - x_{\min}} \]
runs from 0 to 1. Income usually enters through its
logarithm, for the reason of \cref{prop:g12:income-and-wellbeing:doubling}. The choice of bounds and weights is a judgement, not a measurement: it decides how many years of
life are worth a given rise in income, and a different choice can change the ranking. Because the indicators move closely with income per person, as health and satisfaction do,
the rankings by income and by \omterm{def:g12:income-and-wellbeing:composite-indicator}{composite indicators} are often similar; where they differ, the difference is the interesting part.

The table builds one for thirteen countries. Life expectancy is rescaled between 20 and 85 years, the logarithm of income between those of 100 and 75,000 dollars
(capped at 1), and the two are averaged with equal weights; the bounds and the weights are this book's choice. Ranked by income, the United States is first and Japan
fifth; ranked by the composite, Germany and Japan share first place and the United States is fourth. Russia falls from sixth to eighth and Saudi Arabia from third to
fifth, while Nigeria, richer than India, ends below it. Every one of these moves comes from life expectancy, the part that income does not buy at the same rate
everywhere. Put more weight on income and the ranking moves back towards the income ranking; that is the judgement the definition warned of.
% figdata: composite (r_y, r_c, comp)

\begin{omfigure}
\small
\pgfplotstabletypeset[col sep=comma, columns={name,gdppc,le,comp,r_y,r_c},
  columns/name/.style={string type, column name={country}, column type=l},
  columns/gdppc/.style={column name={income (dollars)}, fixed, precision=0, 1000 sep={\,}},
  columns/le/.style={column name={life expectancy}, fixed, precision=1, zerofill},
  columns/comp/.style={column name={composite}, fixed, precision=3, zerofill},
  columns/r_y/.style={column name={rank by income}},
  columns/r_c/.style={column name={rank by composite}},
  every head row/.style={before row=\toprule, after row=\midrule}, every last row/.style={after row=\bottomrule}]
  {figdata/grade-12/15-income-and-wellbeing/out/composite.csv}

\medskip
{\small Thirteen countries in 2019: GDP per person at constant purchasing power, life expectancy at birth, and a composite of the two (equal weights; life
expectancy rescaled between 20 and 85 years, the logarithm of income between those of 100 and 75,000 dollars, capped at 1). Germany and Japan tie at 0.958. Sources: World Bank,
World Development Indicators (SP.DYN.LE00.IN, NY.GDP.PCAP.PP.KD); composite computed in this book.}
\end{omfigure}

\section{What economics measures}\label{sec:g12:income-and-wellbeing:closing}

This book has measured output, prices, wages, trade, debt, inflation and unemployment, and it has tested models against them. These measures are reliable, comparable across
countries and years, and they track a great deal of what people say they value: health, free time and their own satisfaction all rise with income. They do not track
everything, and a measure that leaves something out says nothing about how much it matters. The next volume takes the same tools further, with more mathematics and more data,
and with the same rule: say what a model assumes, what it predicts, how it was tested, and where it breaks.

\section{Exercises}

\begin{exercise}[$\star$]\label{exo:g12:income-and-wellbeing:1}
From the opening chart, give Finland's income and \omterm{def:g12:income-and-wellbeing:life-satisfaction}{life satisfaction} in 2018.
\end{exercise}

\begin{exercise}[$\star$]\label{exo:g12:income-and-wellbeing:2}
From the opening chart, compare Bulgaria (27,691 dollars, 5.4) with Romania (34,829 dollars, 7.4). What does the comparison say about income alone?
\end{exercise}

\begin{exercise}[$\star$]\label{exo:g12:income-and-wellbeing:3}
With a gain of 0.93 points per doubling, how much more satisfied would a country four times richer be? Eight times?
\end{exercise}

\begin{exercise}[$\star$]\label{exo:g12:income-and-wellbeing:4}
German workers worked 2,427 hours a year in 1950 and 1,372 in 2019. By what percentage did hours fall?
\end{exercise}

\begin{exercise}[$\star$]\label{exo:g12:income-and-wellbeing:5}
A country's indices are 0.8 for income, 0.6 for health and 0.7 for schooling. Find the \omterm{def:g12:income-and-wellbeing:composite-indicator}{composite indicator} with equal weights.
\end{exercise}

\begin{exercise}[$\star\star$]\label{exo:g12:income-and-wellbeing:6}
A country at 10,000 dollars has satisfaction 6.0, and one at 40,000 has 7.8. Fit $s = a + b \ln y$ through the two points, and give the gain per doubling.
\end{exercise}

\begin{exercise}[$\star\star$]\label{exo:g12:income-and-wellbeing:7}
Fit the same model through Kosovo (11,188 dollars, 5.9) and Luxembourg (131,734 dollars, 7.5). Compare the gain per doubling with the 0.93 of the full fit. What does this say
about fitting through two points?
\end{exercise}

\begin{exercise}[$\star\star$]\label{exo:g12:income-and-wellbeing:8}
With 0.93 points per doubling, find the gain from a 10\,\% rise in income.
\end{exercise}

\begin{exercise}[$\star\star$]\label{exo:g12:income-and-wellbeing:9}
Country X has an income index of 0.9 and a health index of 0.6; country Y, 0.6 and 0.8. Rank them with equal weights, then with weights of one-third on income and two-thirds
on health.
\end{exercise}

\begin{exercise}[$\star\star$]\label{exo:g12:income-and-wellbeing:10}
From the hours chart, South Korean workers worked 2,749 hours in 1960 and 1,963 in 2019. By how many hours a week did they work less, over 50 weeks?
\end{exercise}

\begin{exercise}[$\star\star$]\label{exo:g12:income-and-wellbeing:11}
A worker earning 30 dalens an hour cuts her hours from 1,800 to 1,600 a year. By how much does her income fall? If she chose the change freely, what is the 200 hours of free
time worth to her, at least?
\end{exercise}

\begin{exercise}[$\star\star\star$]\label{exo:g12:income-and-wellbeing:12}
Give two observations that would contradict the model of satisfaction in the logarithm of income, and say whether any has been found.
\end{exercise}

\begin{exercise}[$\star\star\star$]\label{exo:g12:income-and-wellbeing:13}
Why is the slope of the opening chart not the effect of raising a country's income? Name two reasons.
\end{exercise}

\begin{exercise}[$\star\star\star$]\label{exo:g12:income-and-wellbeing:14}
Workers in country F produce 60,000 dollars a year in 1,400 hours; in country G, 65,000 dollars in 1,800 hours. Value the extra 400 hours of free time in F at G's output per
hour and compare the adjusted figures.
\end{exercise}

\begin{exercise}[$\star\star\star$]\label{exo:g12:income-and-wellbeing:15}
Two countries have the same average income, 40,000 dollars. In A, half the people earn 20,000 and half 60,000; in B, half earn 30,000 and half 50,000. With $b = 1.346$, find
the difference in average satisfaction predicted by the log model, and the income that, given to everyone, would give each country its average satisfaction.
\end{exercise}

\section{Problem: Five Countries, One Ladder}

\begin{problem}[{Weekend problem --- income, life and satisfaction}]
\label{pb:g12:income-and-wellbeing:1}
Five invented countries report their income per person (dollars), annual hours worked per worker, life expectancy (years) and mean \omterm{def:g12:income-and-wellbeing:life-satisfaction}{life satisfaction} (0--10).

\begin{center}
\small
\pgfplotstabletypeset[col sep=comma, columns={country,income,hours,lifeexp,lifesat},
  columns/country/.style={string type, column name={country}, column type=l},
  columns/income/.style={column name={income}, int detect},
  columns/hours/.style={column name={hours}, int detect},
  columns/lifeexp/.style={column name={life expectancy}, int detect},
  columns/lifesat/.style={column name={satisfaction}, fixed, precision=1, zerofill},
  every head row/.style={before row=\toprule, after row=\midrule}, every last row/.style={after row=\bottomrule}]
  {figdata/grade-12/15-income-and-wellbeing/out/five.csv}
\end{center}

\medskip
\noindent\textbf{Part I --- Fitting the model.}
\begin{enumerate}
\item Let $x = \log_2(y/1{,}000)$. Find $x$ for each country.
\item Find the mean of $x$ and the mean satisfaction.
\item Find the least-squares slope of satisfaction on $x$: $\sum (x - \bar x) s / \sum (x - \bar x)^2$. This is the gain per doubling.
\item Write the fitted line.
\item Predict the satisfaction of a sixth country at 128,000 dollars.
\end{enumerate}

\medskip
\noindent\textbf{Part II --- Reading the fit.}
\begin{enumerate}[resume]
\item Find the five residuals.
\item Find the predicted gain from a 10\,\% rise in income.
\item How much income would country C need to reach a satisfaction of 7.0?
\item By how many hours a year does E work less than A?
\item Does the fit show that raising C's income to E's would raise its satisfaction to E's? Explain.
\end{enumerate}

\medskip
\noindent\textbf{Part III --- A composite indicator.}
\begin{enumerate}[resume]
\item The income index is $(\ln y - \ln 1{,}000)/(\ln 128{,}000 - \ln 1{,}000)$. Find it for each country.
\item The life index is $(L - 50)/(90 - 50)$. Find it for each country.
\item Find the \omterm{def:g12:income-and-wellbeing:composite-indicator}{composite indicator}, the average of the two indices.
\item Rank the countries by income, life expectancy, satisfaction and the composite. Do the rankings differ?
\item One more year of life expectancy raises the composite by how much? What rise in income, in per cent, would raise it by the same amount?
\end{enumerate}

\medskip
\noindent\textbf{Part IV --- Judgements in the numbers.}
\begin{enumerate}[resume]
\item By the fitted satisfaction line, what gain in satisfaction goes with that income rise?
\item With weights of one-third on income and two-thirds on life, what income rise would match one more year of life?
\item Which bounds in Part III are choices rather than measurements?
\item What data on these five countries would contradict the log model?
\item State the result of the problem: the satisfaction gain per doubling of income, and the income gain that matches one more year of life expectancy under the equal-weight
  composite.
\end{enumerate}
\end{problem}
